% Part: first-order-logic
% Chapter: beyond
% Section: many-sorted-logic

\documentclass[../../../include/open-logic-section]{subfiles}

\begin{document}

\olfileid{fol}{byd}{msl}

\olsection{منطق چندنوعی}

در منطق مرتبهٔ اول، متغیرها و کمیت‌گذارها روی یک !!{domain} واحد
دامنه می‌گیرند. اما غالباً داشتن چند !!{domain}s (جدا از هم) سودمند است:
برای نمونه، ممکن است بخواهید !!a{domain} از اعداد، !!a{domain} از اشیای
هندسی، !!a{domain} از تابع‌های اعداد به اعداد، !!a{domain} از گروه‌های
آبلی، و مانند این‌ها داشته باشید.

منطق چندنوعی چنین چارچوبی فراهم می‌کند. کار را با فهرستی از «نوع‌ها»
آغاز می‌کنیم---«نوع» یک شیء نشان می‌دهد که آن شیء قرار است در کدام
«!!{domain}» جای داشته باشد. سپس برای هر نوع !!{variable}s و
کمیت‌گذارهایی، و (معمولاً) برای هر نوع یک !!{identity}، داریم. تابع‌ها
و رابطه‌ها نیز بر حسب نوع اشیایی که می‌توانند به‌عنوان آرگومان بپذیرند
«نوع‌بندی» می‌شوند. در دیگر موارد، قواعد معمول منطق مرتبهٔ اول را حفظ
می‌کنیم و نسخه‌های قواعد کمیت‌گذار را برای هر نوع تکرار می‌کنیم.

برای نمونه، در مطالعهٔ روابط بین‌الملل ممکن است زبانی با دو نوع شیء،
شهروندان فرانسوی و شهروندان آلمانی، برگزینیم. ممکن است رابطه‌ای
یک‌موضعی به نام «شراب می‌نوشد» برای اشیای نوع نخست؛ رابطه‌ای یک‌موضعی
دیگر به نام «سوسیس آلمانی می‌خورد» برای اشیای نوع دوم؛ و رابطه‌ای
دوموضعی به نام «یک زوج متأهل چندملیتی تشکیل می‌دهند» داشته باشیم که دو
آرگومان می‌گیرد، به‌گونه‌ای که آرگومان نخست از نوع اول و آرگومان دوم از
نوع دوم باشد. اگر متغیرهای $a$، $b$، $c$ را برای دامنه‌گرفتن روی
شهروندان فرانسوی و $x$، $y$، $z$ را برای دامنه‌گرفتن روی شهروندان
آلمانی به کار ببریم، آنگاه
\[
\lforall[a][\lforall[x][(\Atom{\Obj{MarriedTo}}{a,x} \lif
(\Atom{\Obj{DrinksWine}}{a} \lor \lnot \Atom{\Obj{EatsWurst}}{x}))]]
\]
می‌گوید اگر هر فرانسوی‌ای با یک آلمانی ازدواج کرده باشد، یا آن شخص
فرانسوی شراب می‌نوشد یا آن شخص آلمانی سوسیس آلمانی نمی‌خورد.

منطق چندنوعی را می‌توان به‌شیوه‌ای طبیعی در منطق مرتبهٔ اول نشاند:
همهٔ اشیای !!{domain}s چندنوعی را در یک !!{domain} مرتبهٔ اول گرد
می‌آوریم، با !!{predicate}s یک‌موضعی نوع‌ها را از هم بازمی‌شناسیم، و
کمیت‌گذارها را نسبی می‌کنیم. برای نمونه، زبان مرتبهٔ اول متناظر با مثال
بالا، افزون بر رابطه‌های دیگری که توصیف شدند و پس از حذف الزامات نوعی،
!!{predicate}s یک‌موضعی «$\Obj{German}$» و «$\Obj{French}$» را خواهد
داشت. یک کمیت‌گذار نوع‌دار $\lforall[x][!A]$، که در آن $x$
!!a{variable} از نوع آلمانی است، به صورت
\[
\lforall[x][(\Atom{\Obj{German}}{x} \lif !A)].
\]
ترجمه می‌شود. باید اصول موضوعی‌ای بیفزاییم که جدایی نوع‌ها را تضمین
کنند---برای نمونه، $\lforall[x][\lnot (\Atom{\Obj{German}}{x} \land
  \Atom{\Obj{French}}{x})]$---و نیز اصول موضوعی‌ای که تضمین کنند
«شراب می‌نوشد» تنها دربارهٔ اشیایی برقرار است که محمول
$\Atom{\Obj{French}}{x}$ را برآورده می‌کنند، و مانند آن. با این
قراردادها و اصول موضوع، نشان‌دادن اینکه !!{sentence}s چندنوعی به
!!{sentence}s مرتبهٔ اول و !!{derivation}s چندنوعی به
!!{derivation}s مرتبهٔ اول ترجمه می‌شوند دشوار نیست. همچنین
!!{structure}s چندنوعی به !!{structure}s مرتبهٔ اول متناظر «ترجمه»
می‌شوند؛ و برعکس، هر ساختار مرتبهٔ اولی که اصول جدایی و پوشش نوع‌ها و
قیود نوعی افزوده را برآورده کند، ساختاری چندنوعی به دست می‌دهد. پس برای
منطق چندنوعی نیز قضیهٔ تمامیت داریم.

\end{document}
